\subsection*{先分清旧事实在哪里}

取一条独立的外部付款历史。Alice 的平台付款账户初始有 CNY 10000 分，购买库没有购买。她用 \code{S=(1,purchase.create.external.v2,E1,K1)} 确认购买一份 \code{paper-001}，报价 Q7，数量 1，金额 3000 分。登记会话 E1 处于有效状态。A 共同保存 P1/C1/T1、原输入和固定平台范围内的 \code{pay-C1}。平台 E 建立 X1，余额成为 7000。随后 R 保存成功证据，在实际 C1 中填写 \code{provider\_result} 与完整 \code{payment\_result}，建立 O1 和本地事件 EV1，把 P1 的 \code{business\_state} 推到 \code{paid}，并结束 T1。

设购买数据库的主副本 DB-A 已确认 R，备副本 DB-B 已经持久保存包含原 R 的 WAL，业务状态却只重放到了 A。此刻 DB-B 能读到 P1/C1/T1，P1 显示 \code{awaiting\_payment}，C1 没有成功结果，查询也看不到 O1 和 EV1。原 O1/EV1 的变更依据已经在 DB-B 的持久日志中，平台账本则已经登记 X1，Alice 的平台余额为 7000。

DB-B 接班时先沿已保存的日志重放到原 R，再开放相应恢复入口。原 R 使 O1、EV1、C1 的完整原结果和 T1 的结束事实共同可见，同参请求便从这些行取得原响应。若还要核对平台，恢复者从原 C1 取出平台范围和 \code{pay-C1} 查询 X1。看见 O1 暂时缺失就改用 \code{pay-C2}，会给平台另一项付款身份，余额可能继续降为 4000。

再取一条 DB-B 只保存 A、连原 R 的日志也丢失的历史。C1 与 X1 能确定原付款，却没有保存 R 中实际分配的 O1 和 EV1。假设 EV1 在故障前已经按初次履约身份 \code{initial:O1} 形成 F501，恢复者重新执行分配编号的 R 得到 O2/EV2，下游就可能按 \code{initial:O2} 合法形成 F502。平台没有再扣款，订单、原响应和履约身份却已经改变。恢复者先暂停会重新分配订单或事件身份的 R 重建，并限制首次付款入口，取得独立保存的已提交结果证据，再按原 O1/EV1 重建。第十四章把受理与结果分别保存在独立的权威记录中，具体展开这种恢复。

另一条历史把故障点前移。DB-A 已确认 A，DB-B 连 A 也没有。DB-A 随后失效，DB-B 被提升为新主，入口按 S 查不到 P1。如果这次提升允许舍弃已经确认的 A，原平台映射就从当前购买库消失了。此时重新创建 C2 会与先前可能已发出的 \code{pay-C1} 分开。前十章采用的已确认提交在恢复后持续保留这一前提，已经在这次提升中发生变化。这样的部署要先保存好已确认历史，或进入取得外部证据的受控恢复过程，再开放首次付款。

\CElevenFigure{ch11-original01}{原 R 已持久保存时先重放，原 R 丢失时从独立结果证据恢复原身份}{fig:ch11-two-failovers}

这里要分清三件事：事务是否在原数据库提交，相关日志是否已经保存到接班副本，以及接班副本是否已经把日志重放成可查询的业务状态。DB-B 暂时看不到 O1，可以来自日志尚未到达，也可以来自日志已经持久保存、查询所见状态尚未追上。查询的空结果只给出了当前读取路径的观察，后续动作要结合复制位置、入口权限和原操作仲裁来决定。

\subsection*{保留数据库事务的路径}

在数据库高可用路径中，购买服务继续使用第七至十章的数据结构和事务。A、R 和消费者 Q 的原子提交由各自数据库负责。数据库的复制与恢复机制把这些提交传播给备副本，部署再决定哪些副本有资格接收写入、什么时候确认提交、故障时允许提升谁。应用检查 S、C、EV、I/B 的规则保持，新增的工程问题集中在这些规则依赖哪一份权威历史~\cite{pg}。

以 PostgreSQL 18 的物理流复制为例，数据库通过 WAL 传递恢复所需的写前日志，备库接收、写入、持久刷新，再重放这些记录。WAL 是数据库为故障恢复保存的变更依据，它的内容和本章稍后由应用定义的业务命令日志处于不同抽象层。物理重放在数据库内部恢复页面与事务状态，购买服务通过数据库接口读取已恢复的行。\footnote{WAL 即 Write-Ahead Logging，意为写前日志。PostgreSQL 的物理流复制传递数据库恢复记录。同步提交等待备库刷盘与备库应用在文档中分别说明，部署还要配置同步备库选择~\cite{pg}。}

异步复制允许主库在备库追上以前给客户端确认。因此主库失效时，备库可能缺少已经响应的事务。同步复制可以把确认条件放到指定同步备库的某个阶段。PostgreSQL 的 \code{synchronous\_commit=on} 在配置了同步备库时，等待相应备库持久刷新 WAL；\code{remote\_apply} 进一步等待它重放提交，使事务对该备库查询可见。\code{remote\_write} 等待写入备库操作系统的文件系统，备库操作系统崩溃时的保存条件与持久刷新不同~\cite{pg}。

\CElevenFigure{ch11-wal-stages}{同步提交等待的阶段，决定接班副本已经保存或已经应用什么}{fig:ch11-wal-stages}

这些设置需要连同同步备库选择与故障提升规则一起解释。假设 DB-A 的提交等到了 DB-B 持久刷新，DB-C 却落后。提升 DB-B 可以利用已保存的日志恢复这项提交，提升缺少它的 DB-C 则有另一种历史。界面上启用同步复制以后，故障管理程序还要按实际复制位置选择接班者，并隔离旧主的写入入口，避免旧主恢复后与新主分别受理业务。日志保存位置与主副本资格共同决定此次切换怎样延续已确认提交。

同一事务的数据也要从同一可解释状态读取。R 在一个事务中共同发布 C1 的结果、O1、EV1、P1 的进度与 T1 的结束事实。恢复者如果把 P1 从较新的主库读取、C1 从落后的备库读取，再拼成一个响应，就可能制造出 \code{paid} 配上空支付结果的组合。第九章采用同一读取快照取得固定 \code{acceptance}、P 的当前 \code{progress} 和 C 的原 \code{payment\_result}。副本路由也要把这些读取送入同一可解释状态，才能恢复事务共同发布的关系。

\CElevenFigure{ch11-original02}{事务内共同保存的关系，通过同一读取状态恢复为完整业务响应}{fig:ch11-one-observation}

对恢复入口，可以选择始终读写当前权威主库，让唯一约束和事务仲裁决定首次受理。也可以采用带有明确一致性保证的副本读取接口，并在执行新效果以前回到权威仲裁。备库上的查询结果适合怎样使用，取决于接口承诺的读取新鲜度。一个用于展示的旧状态，可以暂时落后。一个将据以决定首次付款的空结果，则要经过原身份和写入资格的检查。

\subsection*{自己定义状态和命令的路径}

另一条路径从应用开始定义复制单元。把一次合法状态改变写成命令，各副本以相同顺序执行这些命令，并从相同初态得到相同结果。这种组织称为复制状态机。状态机保存当前业务数据，命令描述允许怎样改变它，复制协议负责让有资格执行的副本对命令顺序取得一致~\cite{state-machine-approach}。

例如，三个副本 N1、N2、N3 都保存账户余额、购买、订单、流水、身份绑定和原结果。一条命令携带完整 S、规范化版本和已确认购买输入。副本每次应用命令时，先恢复已有身份，或在新身份上共同产生业务变化与结果。这里共同发布的边界由应用状态机及其持久化路径构造，开发者要负责输入解释、编号分配、快照格式和升级兼容。

本章用一个封闭的本地余额模型推导这条路径。它沿用第一部分的动作 \code{purchase.\allowbreak create.\allowbreak v1}，与第二部分在运行的 \code{purchase.\allowbreak create.\allowbreak external.\allowbreak v2} 分开。每条历史都明确给出初态，余额只含购买支出，金额来自预置、可验证的报价 Q7。把这个模型用于讲复制，可以直接看到重复命令造成的两次减款。现有外部支付链路继续由原数据库处理 A/R，由平台处理 E。

\CElevenFigure{ch11-original03}{数据库复制传播数据库提交，独立状态机复制传播应用命令}{fig:ch11-replication-paths}

若把外部支付链路迁入自建复制状态机，需要先定义相当于 A/R 的确定性状态迁移，再定义状态机外部的发送者怎样读取持久待办、按 C1 调用平台，以及怎样将可信证据提交回状态机。三个副本在应用同一条命令时各向平台发一次扣款，会产生三次外部调用。平台能否把这三次合并，依赖原范围和 \code{pay-C1} 的规则。应用命令先保存待办和固定付款身份，外部发送者再据此工作，才把复制前缀接回原支付协议。

同样，履约库也可以独立采用数据库高可用或复制状态机。Q 要共同保留 I、按 B 仲裁的真实 F、原接收结果和回报待办。F1 的实际编号 501、平台通知 PEV 与本地事件 EV1 的区别，都来自对应事务域中的业务事实。购买库、履约库和平台分别保存自己的提交，再沿 C、EV1 与 I/B 传递证据、恢复后续工作。

\CElevenFigure{ch11-domains}{复制在各自提交域内延续事实，恢复关系跨域传递固定身份与证据}{fig:ch11-domains}
\FloatBarrier
